\documentclass[dvipdfmx]{jarticle}
\hoffset=-25mm \voffset=-25mm
\setlength{\textwidth}{16cm}
\setlength{\textheight}{23cm}
\usepackage{url}
\usepackage{tikz}
\usetikzlibrary{calc,shapes}
\newcounter{coffee}
\newcounter{bean}
\begin{document}
\begin{flushright}
  \texttt{2026/10/16}
\end{flushright}
\begin{center}{\bf\Large
    2026年度　コンパイラ演習課題(3)
  }\end{center}
\begin{itemize}
  \item TACTの課題の添付ファイルとしてPDFファイルで提出すること．
  \item 手書きレポートをPDFとして提出方法は講義webサイトを
        参照\footnote{\url{https://sites.google.com/sqlab.jp/26compiler/how-to-make-report}}すること．
  \item 提出期限: 10月27日(火) 23:55(JST)
\end{itemize}
\begin{list}{[\arabic{coffee}]}{\usecounter{coffee}}
    \item 以下の文法について答えよ．
        \[\langle\{S,E\},\{id,\texttt{(},\texttt{)},
        \left\{\begin{array}{clcl}
            0: & S\rightarrow E\ \# & 2: & E\rightarrow id\ \texttt{(}\ E\ \texttt{)}\}\\
            1: & E\rightarrow id & 3: & E\rightarrow E\ \texttt{+}\ E\\
        \end{array}\right\},S
        \rangle\]
        \begin{list}{(\alph{bean})}{\usecounter{bean}}
            \item この文法はLR(0)文法か？(LR(0)オートマトンに基づいて先読み文字なしで
            PDTの操作を決定して還元が可能か？)
            \item この文法はSLR(0)文法か？
            \item この文法はLR(1)文法か？
        \end{list}
    \item 以下の文法はLL(1)解析可能ではあるがLALR(1)解析可能でないことを示せ．
        \[\langle\{S,A,E,F,X\},\{\texttt{(},\texttt{)},\texttt{[},\texttt{]},\left\{
            \begin{array}{clcl}
                1: & S\rightarrow \verb|(|X & 5: & X\rightarrow F\verb|]| \\
                2: & S\rightarrow E\verb|]| & 6: & E \rightarrow A        \\
                3: & S\rightarrow F\verb|)| & 7: & F\rightarrow A         \\
                4: & X\rightarrow E\verb|)| & 8: & A\rightarrow
            \end{array}
        \right\},S\rangle\]
　  \item 以下の文法はLALR(1)解析可能であるがLL(1)解析可能でないことを示せ．
        \[\langle\{S,A,B\},\{\texttt{a},\texttt{b},\texttt{c},\texttt{d}\},\left\{
            \begin{array}{clcl}
                1: & S\rightarrow A\verb|a|         & 4: & S\rightarrow \verb|b|B\verb|a| \\
                2: & S\rightarrow \verb|b|A\verb|c| & 5: & A\rightarrow\verb|d|           \\
                3: & S\rightarrow B\verb|c|         & 6: & B\rightarrow\verb|d|
            \end{array}\right\},S\rangle\]
    \item 以下の文法はLALR(1)解析可能であるがSLR(1)解析可能でないことを示せ．
             \[\langle\{S,X,M\{,\{\texttt{a},\texttt{b},\texttt{c},\texttt{d}\},\left\{
            \begin{array}{clcl}
                0: & S\rightarrow X         & 3: & X\rightarrow \verb|d|\verb|c| \\
                1: & X\rightarrow M \verb|a| & 4: & X\rightarrow\verb|b|\verb|d|\verb|a|\\
                2: & X\rightarrow \verb|b|M\verb|c|         & 5: & X\rightarrow\verb|d|
            \end{array}\right\},S\rangle\]
    \item 以下の文法について，LALR(1)構文解析表を作成すると複数の動作をもつエントリが存在する．
            \[\langle\{S,E\},\{\texttt{while},\texttt{do},\texttt{:=},\texttt{+},id\},
            \left\{\begin{array}{clcl}
                0: & S\rightarrow E\ \# & 3: & E\rightarrow E\ \texttt{+}\ E\\
                1: & E\rightarrow \texttt{while}\ E\ \texttt{do}\ E & 4: & E\rightarrow id\\
                2: & E\rightarrow id\texttt{:=}E & & \\
            \end{array}\right\},S\rangle\]
            この複数の動作がある状態について説明し，一般的な構文解析木を構成するためには，
            どちらの動作を優先すべきであるかについて解答せよ．
\end{list}
\end{document}
